\section*{Chapter \ref{ch:g4:geometry} --- Lines and Polygons}

\begin{solution}{exo:g4:geometry:1}
Constructions with \cref{met:g4:geometry:draw}; the little square
marks the right angle between $d$ and the perpendicular.
\end{solution}

\begin{solution}{exo:g4:geometry:2}
E: the three horizontal strokes are parallel; each is perpendicular
to the vertical one. H: the two vertical strokes are parallel, the
crossbar is perpendicular to both. N: the two vertical strokes are
parallel; the diagonal is neither. T: the two strokes are
perpendicular. Z: the top and bottom strokes are parallel; the
diagonal is neither.
\end{solution}

\begin{solution}{exo:g4:geometry:3}
Hexagon; quadrilateral; triangle; octagon; pentagon. A polygon has
as many vertices as sides: $6$, $4$, $3$, $8$, $5$.
\end{solution}

\begin{solution}{exo:g4:geometry:4}
A pentagon has $5$ sides. From $A$, the non-neighbors are $C$ and
$D$: two diagonals, $[AC]$ and $[AD]$.
\end{solution}

\begin{solution}{exo:g4:geometry:5}
$OA = 4$ cm (radius); $OM = 4$ cm (every point of the circle);
$AB = 8$ cm (a diameter is twice the radius).
\end{solution}

\begin{solution}{exo:g4:geometry:6}
Open the compass along one side, then carry that opening to the two
other sides: if the spike-and-pencil span fits a second side
exactly, those two sides are equal.
\end{solution}

\begin{solution}{exo:g4:geometry:7}
The path closes into a polygon with $5$ sides --- a pentagon (not a
regular one!).
\end{solution}

\begin{solution}{exo:g4:geometry:8}
``Perpendicular lines always cross'': true --- they cross at the
right angle itself.

``Parallel lines never cross'': true --- that is their definition.

``A line can be parallel to one line and perpendicular to
another'': true --- on grid paper, a horizontal line is parallel to
every horizontal line and perpendicular to every vertical one.
\end{solution}

\begin{solution}{exo:g4:geometry:9}
The lines $d$ and $f$ look parallel --- and they always are: two
lines perpendicular to the same line are parallel
(\cref{prop:g6:lines:facts} makes this official).
\end{solution}

\begin{solution}{exo:g4:geometry:10}
One correct program: ``1. Draw a square $ABCD$ of side $4$ cm.
2. Mark the point $O$ where its diagonals cross. 3. Draw the circle
of center $O$ and radius $2$ cm.'' The circle touches each side at
its middle.
\end{solution}

\begin{solution}{exo:g4:geometry:11}
Quadrilateral: $2$ diagonals. Pentagon: $5$. Hexagon: $9$. Each new
vertex adds more and more diagonals --- the counts grow faster and
faster ($+3$, then $+4$, \dots).
\end{solution}
